%----------------------------------------------------------------------------
\chapter{Conclusion and future work}
\label{chap:Conclusion}
%----------------------------------------------------------------------------

%----------------------------------------------------------------------------
\section{Summary}
\label{sec:Summary}
%----------------------------------------------------------------------------
\begin{itemize}
    \item This thesis presented a controlled comparison of XPBD and VBD, one of the most dominant solvers for real-time deformable simulation and a relatively new approach.
    \item The contribution is threefold:
    \begin{itemize}
        \item a shared Unity codebase implementing both solvers under matched material parameters and compute budget;
        \item a Newton ground-truth reference, enabling reporting of solution error rather than only relative performance;
    \end{itemize}
\end{itemize}
\emph{TODO: once results are finalised, a single paragraph summary of
\autoref{sec:Verdicts} goes here: how each solver behaves in the scenarios the
primal/dual lens flags, what the equal-wall-clock frontier looks like, and the
scenario-based recommendation for practitioners.}

%----------------------------------------------------------------------------
\section{Future work}
\label{sec:FutureWork}
%----------------------------------------------------------------------------

\subsection{Multithreaded and GPU comparison}
\label{ssec:FWParallel}
\begin{itemize}
    \item The most consequential extension is to repeat the experiments on a multithreaded CPU and on the GPU, directly lifting the single-threaded limitation of \autoref{ssec:ThreatSingleThread}.
    \item Both solvers admit parallelisation via graph colouring, but the structures differ in granularity (\autoref{sec:Parallelisation}): VBD's per-vertex blocks should colour more cleanly than XPBD's per-constraint blocks.
    \item The expected outcome is therefore that parallelisation widens VBD's lead; quantifying this gap is the natural next study.
\end{itemize}

\subsection{Extend research to other domains of dynamics}
\label{ssec:FWOtherDynamics}
\begin{itemize}
    \item Beside mass-spring systems, in theory both solvers could also handle other cloth representations, soft bodies, and rigid bodies as well.
    \item It would be interesting to see how the two solvers perform on these other types of dynamics, and whether the same conclusions hold.
\end{itemize}

\subsection{Using a different cloth model}
\label{ssec:FWFEM}
\begin{itemize}
    \item The mass-spring discretisation is the source of the bias discussed in \autoref{ssec:ThreatMassSpring}, which makes the cloth model the most informative ingredient to vary next.
    \item The cheapest change is to replace my 12-neighbour-per-vertex model with a 4- or 8-neighbour one; this narrows the range of representable materials but would reveal how sensitive the comparison is to mesh connectivity.
    \item The more substantial change is to swap the mass-spring energy for an energy-based model such as Saint-Venant--Kirchhoff (StVK), the model the original VBD paper uses; I would expect it to favour XPBD far less.
    \item A particularly informative variant would pair the models asymmetrically --- mass-spring XPBD against StVK VBD --- letting each solver run on the discretisation that suits it best.
\end{itemize}

\subsection{AVBD and post-2024 variants}
\label{ssec:FWAVBD}
\begin{itemize}
    \item Augmented VBD~\cite{Giles2025} reportedly handles wide stiffness ratios and high stiffness substantially better than canonical VBD.
    \item Re-running the stiffness-ratio experiments with AVBD would resolve the question left open in \autoref{ssec:ThreatAVBD}: whether the gap I observe is intrinsic to the primal approach or an artefact of the canonical 2024 implementation.
\end{itemize}
